\chapter{Channel（渠道）}

\vspace{0.5em}
\noindent\textit{上一章我们认识了 Agent 智能体——OpenClaw 的“大脑”。但智能体需要与外部世界沟通，
本章将介绍 Channel（渠道）——连接各类聊天平台的适配器。我们将了解内置渠道、登录方式、
访问控制策略以及路由机制。}
\vspace{0.5em}

\section{概述}

Channel（渠道）是 OpenClaw 与聊天平台之间的适配器。每个渠道负责：

\begin{itemize}
  \item 接收平台入站消息并转换为统一格式
  \item 将智能体回复发送回对应平台
  \item 处理平台特有功能（附件、回应、线程等）
\end{itemize}

配置路径：\texttt{channels.<name>.*}

\section{内置渠道}

\begin{center}
\begin{tabular}{lll}
\toprule
\textbf{渠道 ID} & \textbf{平台} & \textbf{登录方式} \\
\midrule
\texttt{whatsapp} & WhatsApp & QR 码配对 \\
\texttt{telegram} & Telegram & Bot Token \\
\texttt{discord} & Discord & Bot Token \\
\texttt{slack} & Slack & OAuth/Token \\
\texttt{signal} & Signal & 配对 \\
\texttt{imessage} & iMessage（旧版） & macOS \\
\texttt{bluebubbles} & iMessage（推荐） & BlueBubbles 服务器 \\
\texttt{webchat} & Web 聊天 & 内置 \\
\texttt{googlechat} & Google Chat & Service Account \\
\bottomrule
\end{tabular}
\end{center}

插件可添加更多渠道（如 Mattermost、Zalo、WeChat 等）。

\section{渠道登录}

\begin{lstlisting}[language=bash]
# Interactive login (all channels)
openclaw channels login

# Channel-specific
openclaw channels login --channel telegram
openclaw channels login --channel discord

# Check status
openclaw channels status --probe
\end{lstlisting}

\section{访问控制}

每个渠道支持 \textbf{私信策略}和\textbf{群组策略}：

\subsection{私信策略（dmPolicy）}

\begin{itemize}
  \item \texttt{pairing}（默认）：未知发送者收到配对码，审批后方可通信
  \item \texttt{allowlist}：仅允许名单内的用户
  \item \texttt{open}：任何人都可以私信
  \item \texttt{disabled}：禁用私信
\end{itemize}

\subsection{群组策略（groupPolicy）}

\begin{itemize}
  \item \texttt{allowlist}（默认）：仅允许名单内的群组
  \item \texttt{open}：所有群组
  \item \texttt{disabled}：禁用群组
\end{itemize}

\subsection{配置示例}

\begin{lstlisting}[language=json]
{
  channels: {
    whatsapp: {
      dmPolicy: "pairing",
      allowFrom: ["+15555550123"],
      groups: {
        "*": { requireMention: true }
      }
    },
    telegram: {
      enabled: true,
      dmPolicy: "allowlist",
      allowFrom: ["123456789"]
    }
  }
}
\end{lstlisting}

\section{路由机制}

OpenClaw 将回复\textbf{路由回消息来源的渠道}，路由是确定性的：

\begin{itemize}
  \item 模型\textbf{不会}选择渠道
  \item 回复始终返回到入站消息的来源
  \item 多渠道场景通过 Binding 规则路由到不同 Agent
\end{itemize}

\subsection{会话键格式}

\begin{itemize}
  \item 私信：\texttt{agent:<agentId>:<mainKey>}
  \item 群组：\texttt{agent:<agentId>:<channel>:group:<groupId>}
  \item 线程：在基础键后追加 \texttt{:thread:<threadId>}
\end{itemize}

\section{多账户}

支持多账户的渠道（如 WhatsApp）使用 \texttt{accountId} 标识每个登录实例。
不同 \texttt{accountId} 可路由到不同 Agent，实现一台服务器托管多个号码。

\section{本章小结}

本章介绍了 Channel 渠道的核心机制：

\begin{itemize}
  \item Channel 是 OpenClaw 与外部聊天平台之间的\textbf{协议翻译器}
  \item 内置支持 WhatsApp、Telegram、Discord 等 9 个平台，还可通过插件扩展
  \item 提供精细的\textbf{访问控制}（私信策略 + 群组策略）
  \item 路由是\textbf{确定性的}，回复始终返回入站消息的来源渠道
\end{itemize}

\noindent 渠道解决了“消息从哪里来”的问题，但在多智能体场景下，如何决定消息“发给谁”？下一章我们将讲解 Binding（路由绑定）机制。
